% PASS20260921 twin Z6 single-node difference.
\subsection*{The structural and holonomy $Z_6$ classes differ by one mark-two coweight phase}

In the common affine-$E_8$ Kac gauge, let
\[
S=(2,0,0,0,1,0,0,0,0)
\]
denote the structural FI $\times$ matter-parity class and
\[
H=(0,0,0,0,1,0,0,1,0)
\]
the flagship holonomy class.  Only the eight finite Kac labels determine root
phases.  Their quotient is therefore the single finite mark-two node,
normalized to exact order six as
\[
\boxed{R=(4,0,0,0,0,0,0,1,0)}.
\]
The full $240$-root census proves
\[
\boxed{\deg_H(r)=\deg_S(r)+\deg_R(r)\pmod6}
\]
for every root, hence in the common maximal torus
\[
\boxed{H=S\,R}.
\]

The three eigenspace fingerprints are
\[
d(S)=(54,48,33,32,33,48),
\]
\[
\boxed{d(R)=(92,64,14,0,14,64)},
\]
and
\[
d(H)=(44,40,38,48,38,40).
\]
The difference element has power ladder
\[
R:\ D_7+\mathfrak u_1,qquad
R^2:\ D_7+\mathfrak u_1,qquad
R^3:\ D_8,
\]
with fixed dimensions $92,92,120$, and adjoint trace powers
\[
(248,142,14,-8,14,142).
\]
Thus the two physicalized order-six structures share the same mark-four Kac
component; one additional mark-two fundamental-coweight phase accounts for
their complete root-space difference.

\paragraph{Scope firewall.}
This is an exact factorization of inner automorphisms in the common complex
$E_8$ Cartan gauge.  It does not imply that the heterotic FI dynamics
generates the compactification holonomy and does not solve D/F flatness.

The executable witness is
\texttt{analysis/w33\_twin\_z6\_single\_node\_difference.py}.
